\documentclass{subfiles}
\begin{document}
    As for the case of \lstinline{RB-Insert}, 
        the \lstinline{RB-Delete} procedure is also based on the BST one,
        with few differences: the main one being the call to \lstinline{RB-Delete-Fixup}.

        % TODO: write RB-Delete
    
    Once again, we mainly focus on \lstinline{RB-Delete-Fixup} procedure 
        as this does the heavy lift. Consider \Cref{RB-Delete-Fixup}'s code: 
        it's evident that the properties that can fail are property 
        \ref{RB-trees:const1}, \ref{RB-trees:const2}, \ref{RB-trees:const4} and 
        \ref{RB-trees:const5} are violated. We can correct the latter,
        by imposing the following: when a black node \(y\) is removed or moved, 
        the node \(x\) that takes its place gets an ``extra'' black.
        Thus, under this interpretation, \lstinline{RB-Delete-Fixup} must correct 
        properties \ref{RB-trees:const1}, \ref{RB-trees:const2} and \ref{RB-trees:const4}.

    We focus our discussion on proving that the procedure corrects property \ref{RB-trees:const1},
        as the other follow directly. We have eight possible cases, in pairs symmetrical.
        Hence we focus our attention on case 1-4. In our treating of \lstinline{RB-Delete-Fixup},
        \(x\) always points to the node with the ``extra'' black.

    \begin{description}
        \item [Case 1:] \(x\)'s sibling \(w\) is red;
            therefore, if we swap the color of \(x.p\) and \(w\) and rotate to the left on 
            \(w\). In doing so, one of Case 2, Case 3 or Case 4 will fix the tree. 

        \item [Case 2:] \(w\) and its children are black;
            we remove a red from both \(x\) and \(w\) and add an extra black to \(x.p\) 
            to conpensate.

        \item [Case 3:] \(w\) is black, its left child is red and is right child is black. 
            We can simply switch the color of \(w\) and is left child and do a right rotation on 
            \(w\), thus folling in Case 4.

        \item [Case 4:] \(w\) is black and is right child is red;
            we can apply some color change and rotate to the right on \(x.p\),
            without any reb-black property being violated.
    \end{description}

    In terms of time complexity, \lstinline{RB-Delete} is easely seen to be \(\bigO{\log n}\).
    

\end{document}
